\scanpage{157}
\mspage{15}
mais ne sont pas impliquées par celles-ci.

Contr'exemple, la catégorie 0-connexe
\diagram{\begin{tikzcd}[ampersand replacement=\&,column sep=small]
 0\arrow[r,bend left,"\alpha"]\arrow[r,bend right,"\beta"'] \& 1
\end{tikzcd}}
, avec $\alpha\ne\beta$. \unsure{On a|Car.} \deleted{\unsure{$D=$}} $X_{/1}\simeq$
\diagram{\begin{tikzcd}[ampersand replacement=\&,row sep=0pt,column sep=small]
 \& \unsure{0|0'}\arrow[dr] \& {}\\
 \deleted{\unsure{0}}\arrow[ur,dash,"\deleted{\unsure{\alpha}}"]\arrow[dr,dash,"\deleted{\unsure{\beta}}"'] \& {} \& 1\\
 \&0'\arrow[ur] \& {}
\end{tikzcd}}
, qui n'est pas irréductible, puisque $\unsure{D|\Phi}\smallsetminus\{1\}$ \inserted{$\simeq\unsure{2|e}$} est un ouvert \unsure{disconnexe|déconnexe}.

J'aimerais une condition ``toposique'' sur $X$, équivalente à \deleted{\unsure{$I$}} $X$ ``pseudofiltrante à gauche'' i.e. les composantes connexes filtrantes ; i.e. $\mathrm{PS}\underline{1}^{\mathrm{opp}}$ (condition \unsure{c'|b'}\unsure{)} ci-dessus) $+$ $\mathrm{PS}2^{\mathrm{opp}}$ (condition de la double flèche), et qui de plus implique que tout objet \inserted{0-}connexe de $X^\wedge$ (i.e. tt objet $X'$ de $\mathrm{Cat}$ étale sur $X$ et 0-connexe) soit $W_\infty$-asphérique. Et je songe à la notion des ``\unsure{totally asphericats|totales asphéricités}'' de PS (\unsure{Pursuing Stacks}).

\uline{Définition} \unsure{a)} 2.2. Soit $X$ une catégorie. \deleted{\unsure{a}}
% slanted
\margin{Cf tableau récapitulatif pages 97, 98 et surtout 99}
On dit que $X$ est \uline{faiblement filtrante}* \deleted{(ou \unsure{fbt.} filtrante} \uline{à gauche} \inserted{ou \uline{irréductible}}, \deleted{ou \unsure{quand}} si $X\ne\varnothing$, et si pour $a,b\in\mathrm{Ob}\,X$ $\exists c$ et des flèches
\diagram{\begin{tikzcd}[ampersand replacement=\&,row sep=0pt,column sep=tiny]
 \& a\\ c\arrow[ur]\arrow[dr] \& {}\\ \& b
\end{tikzcd}}
, i.e. $\deleted{\unsure{X}} X_{/(a)}\cap\deleted{\unsure{X}}X_{/[b]}\ne\varnothing$.
% slanted
\margin{dire plutôt \uline{pseudo-filtrante} \unsure{(puisque)} \unsure{réserver|c'est une} \unsure{appellation} loc. fbt filtrante'' condition PS $\underline{1}'_{\mathrm g}$}
(\uline{NB} $X$ est ps. filtr. \inserted{à g.} ssi $\mathrm{Ord}(X)$ \inserted{(i.e. est filtr. à g)} l'est.)
\uline{ReNB} $X$ fbt filtrante $\Longrightarrow X$ 0-connexe

b) On dit que $X$ est \uline{loc. faiblement filtrante} \deleted{\unsure{si}} ou
\inserted{\uline{ou plus simpl. $\longrightarrow$ loc. filtrante à gauche}} \deleted{les catégories} \uline{à gauche}, \deleted{\unsure{ou ps.}} si les catégories $X_{/x}$ ($x\in\mathrm{Ob}\,X$) sont fbt filtrantes à gauche, i.e. si $X$ satisfait la condition : $\forall$ diagramme
\diagram{\begin{tikzcd}[ampersand replacement=\&,row sep=0pt,column sep=small]
 a\arrow[dr] \& {}\\{} \&x,\\ b\arrow[ur] \&{}
\end{tikzcd}}
$\exists c$ de $X$ et des flèches de $c$ dans $a,b$ rendant \uuline{commutatif} le carré
% slanted
\margin{condition PS $1_{\mathrm g}$}
\diagram{\begin{tikzcd}[ampersand replacement=\&,row sep=0pt,column sep=small]
 \& a\arrow[dr] \& {}\\ c\arrow[ur,dotted]\arrow[dr,dotted] \&{} \&x\\ \& b\arrow[ur] \&{}
\end{tikzcd}}
[\uline{NB} Cette condition \deleted{n'est pas} implique que $\mathrm{Ord}\,X$ soit loc fbt filtrante, mais n'équivaut pas à cette dernière. Exemple $X={}$%
\diagram{\begin{tikzcd}[ampersand replacement=\&,column sep=small]
 0\arrow[r,bend left,"\alpha"]\arrow[r,bend right,"\beta"'] \&1
\end{tikzcd}}
, $X_{/1}={}$%
\diagram{\begin{tikzcd}[ampersand replacement=\&,row sep=0pt,column sep=small]
 0'\arrow[dr] \&{}\\ {} \&1\\ 0''\arrow[ur] \&{}
\end{tikzcd}}
n'est pas \unsure{faib.} filtr. à g.]
% archive sheet number: 157
